\chapter{Results Tab}
\label{chap:resultats}

The \menu{Results} tab presents simulation results
as a \textbf{configurable grid of analysis
cells}. Each cell displays an independent plot type
(polar, pressure distribution, boundary layer profile,
etc.).

\screenshot{10_tab_results_empty.png}{1.0}{%
    \menu{Results} tab before any simulation. The grid
    defaults to a $2\times2$ layout, i.e.
    four empty cells.}

\screenshot{18_tab_results.png}{1.0}{%
    \menu{Results} tab after simulating the current airfoil
    (blue), the reference (red) and the airfoil with flap (green).
    The four cells here display $C_L(\alpha)$, the
    $C_L(C_D)$ polar, the lift-to-drag ratio $C_L/C_D$ as a function of $C_L$ and the
    $C_M(C_L)$ moment.}

\section{Grid configuration}

The grid is configurable via the
\menu{View $\rightarrow$ Layout} menu. Five preconfigured
layouts are available: $1\times1$, $1\times2$,
$2\times2$ (default), $2\times3$, $3\times3$.

\begin{noteinfo}
When the layout is changed, the analyses already
selected in existing cells are kept.
New cells are initialized with the default
analyses.
\end{noteinfo}

\section{Analysis cell}

Each cell features a \textbf{control bar} and a
\textbf{matplotlib canvas}.

\subsection{Control bar}

\begin{tabularx}{\linewidth}{l X}
    \toprule
    \textbf{Element} & \textbf{Role} \\
    \midrule
    Drop-down list (analysis)
        & Choice of the analysis type to display in the cell
          (see \emph{Analysis types}, below). \\
    \menu{C} checkbox
        & Shows / hides the current airfoil curves (blue). \\
    \menu{R} checkbox
        & Shows / hides the reference airfoil curves (red). \\
    \menu{F} checkbox
        & Shows / hides the airfoil with flap curves (green). \\
    \menu{Reynolds} list
        & For analyses that depend on a particular Reynolds
          ($-C_p(x)$, Cp~+~Airfoil, airfoil~+~CL), selects the
          Reynolds to display. Hidden if not relevant. \\
    \menu{Alpha} list
        & For analyses that depend on an angle of attack
          (Cp~+~Airfoil, airfoil~+~CL and the boundary layer
          quantities), selects the angle of attack $\alpha$ to
          display. Hidden if not relevant. \\
    Coordinates
        & Real-time display of the cursor coordinates
          in the plot frame. \\
    \bottomrule
\end{tabularx}

\subsection{Mouse interactions}

The interactions available on the matplotlib canvas are
identical to those of the \menu{Airfoils} tab:

\begin{itemize}
    \item Left click + drag: rectangle zoom;
    \item Wheel: zoom centered on the cursor;
    \item Middle click: pan;
    \item Left click on the legend: hides / shows the
        corresponding curve.
\end{itemize}

\subsection{Interactive and scrollable legend}

When a cell contains many curves (for example
several Reynolds combined with several airfoils), the legend can
become too long to be displayed entirely. A \textbf{scroll
bar} appears automatically beyond 15~entries.

\begin{astuce}
To quickly compare two curves, it is more efficient
to \textbf{hide} the interfering curves by clicking on their
entry in the legend, rather than zooming or adding
an extra cell.
\end{astuce}

\section{Available analysis types}

\subsection{Polars (as a function of Alpha)}

\begin{tabularx}{\linewidth}{l X}
    \toprule
    \textbf{Analysis} & \textbf{Description} \\
    \midrule
    $C_L(\alpha)$  & Lift coefficient as a function of
                     the angle of attack. One curve per Reynolds.\\
    $C_D(\alpha)$  & Drag coefficient as a function of
                     the angle of attack.\\
    $C_M(\alpha)$  & Pitching moment coefficient
                     (quarter-chord reference).\\
    $C_L/C_D$~$(\alpha)$ & \emph{Lift-to-drag ratio} as a function of
                     the angle of attack. Allows the
                     angle of attack of maximum lift-to-drag ratio to be determined.\\
    $C_L/C_D$~$(C_L)$ & \emph{Lift-to-drag ratio} as a function of the lift
                     coefficient. Convenient for reading the lift-to-drag ratio at
                     the targeted cruise $C_L$.\\
    $C_M(C_L)$     & Moment coefficient as a function of the
                     lift. The slope indicates stability;
                     the effect of a flap is very readable here.\\
    $C_D(C_L)$ \emph{(Lilienthal polar)}
                   & Classic polar: drag on the x-axis
                     vs lift on the y-axis. The tangent
                     from the origin gives the maximum lift-to-drag ratio.\\
    Transition top, bot
                   & Position $x/c$ of the boundary layer
                     transition on the upper / lower surface, as a function
                     of $\alpha$. Indicates the effective
                     roughness and laminar separation.\\
    \bottomrule
\end{tabularx}

\subsection{Distributions along the airfoil}

\begin{tabularx}{\linewidth}{l X}
    \toprule
    \textbf{Analysis} & \textbf{Description} \\
    \midrule
    $-C_p(x)$ & Distribution of the pressure coefficient along
               the chord, at a given Reynolds and angle of attack.
               Convention: $-C_p$ upward (upper surface on top). \\
    Cp + Airfoil & Combined \emph{XFoil}-style plot: $-C_p$ on top and
               the airfoil redrawn below, for a (Reynolds, angle of
               attack) pair. An inset summarizes the coefficients
               ($C_L$, $C_M$, $C_D$, lift-to-drag ratio). For a fuller
               dedicated view (inviscid $C_p$ and boundary layer), see
               chapter~\ref{chap:cp}. \\
    Airfoil + CL & Combined visualization: airfoil rotated by
               the angle of attack + plot of the boundary layer
               edge (displacement thickness $\delta^*$).\\
    \bottomrule
\end{tabularx}

\subsection{Boundary layer}

Five boundary layer quantities can be plotted as a function
of the curvilinear abscissa $s$ or the Cartesian abscissa $x$. They
are evaluated at the \textbf{selected angle of attack} (\menu{Alpha}
list) — the boundary layer being computed for each $\alpha$ — and each
curve compares the available Reynolds numbers:

\begin{tabularx}{\linewidth}{l X}
    \toprule
    \textbf{Quantity} & \textbf{Description} \\
    \midrule
    $U_e(s)$, $U_e(x)$
        & Velocity at the edge of the boundary layer,
          normalized by the upstream free-stream velocity.\\
    $\delta^*(s)$, $\delta^*(x)$
        & Displacement thickness.\\
    $\theta(s)$, $\theta(x)$
        & Momentum thickness.\\
    $C_f(s)$, $C_f(x)$
        & Local skin friction coefficient.\\
    $H(s)$, $H(x)$
        & Shape factor $H = \delta^*/\theta$ ($H \approx 2{,}6$
          at the transition point, $H \approx 4$ at separation).\\
    \bottomrule
\end{tabularx}

\section{Stale results banner}

When the user returns to the \menu{Airfoils} tab and modifies
an airfoil after simulation (moving a control point,
changing the degree, conversion\dots), a yellow banner appears
at the top of the \menu{Results} tab:

\begin{quote}
\textcolor{orange!75!black}{$\triangle$
\emph{The current airfoil has been modified since the last
simulation. Re-run to update.}}
\end{quote}

This banner disappears automatically at the next simulation.

\begin{important}
The results displayed in the cells are \emph{not}
recomputed automatically after an airfoil is modified.
The banner is a simple visual indicator: it is up
to the user to re-run the simulation from the
\menu{XFoil Settings} tab.
\end{important}
